\documentclass[10pt,a4paper]{article}
\usepackage[margin=2cm]{geometry}
\usepackage{amsmath,amssymb}
\usepackage{booktabs}
\usepackage{graphicx}
\usepackage{xcolor}
\usepackage{hyperref}
\usepackage[expansion=false]{microtype}
\usepackage{caption}
\usepackage{subcaption}
\usepackage[T1]{fontenc}

\hypersetup{colorlinks=true,linkcolor=blue!60!black,citecolor=green!50!black}

\definecolor{warn}{RGB}{192,57,43}
\definecolor{good}{RGB}{39,174,96}

\title{\textbf{EB-JEPA for EEG Time Series}\\
  \large Self-Supervised Abnormality Detection on TUAB\\[4pt]
  \normalsize Hackathon Technical Report}
\author{Yani Hammache, Théo Basseras, Titouan Vasnier, Timothy Courtois}
\date{June 2026}

\begin{document}
\maketitle
\begin{abstract}
We apply the Energy-Based JEPA (EB-JEPA) framework to EEG time-series, training a
self-supervised encoder on the TUH Abnormal EEG Corpus (TUAB) with no labels.
A frozen linear probe on the learned representations reaches
\textbf{BACC 0.775} per-window (fair comparison to the literature)
and \textbf{BACC 0.836} per-recording (mean-pool of 16 windows per patient).
We compare four encoder architectures and two SSL objectives, identify a
VICReg coefficient bug (invariance underweighted by 25$\times$) as the primary
cause of EEGPT encoder collapse, and flag a systematic evaluation protocol
mismatch that initially caused us to overstate our results against published baselines.
\end{abstract}

\section{Task and Dataset}
\textbf{TUAB} (TUH Abnormal EEG Corpus v3.0.1): binary normal/abnormal classification
on scalp EEG. 19 channels, 200\,Hz, 10\,s windows, z-scored per channel.
Patient-disjoint split: $n_{\text{train}}=2717$ recordings / $n_{\text{eval}}=276$ recordings
(45.6\% abnormal). SSL pretraining uses only the train split with no labels.
Evaluation metric: \textbf{Balanced Accuracy (BACC)} and AUROC.

\section{Our Contribution}

\paragraph{What is provided (EB-JEPA framework).}
The two-view SSL training loop, dataset loader, data augmentations, VICReg and SIGReg
(BCS/Epps-Pulley) loss implementations, and the Projector MLP are provided by the
competition framework.

\paragraph{What we implemented and contributed.}
\begin{enumerate}
  \item \textbf{Four encoder architectures} (\S\ref{sec:encoders}): strided Conv1D (baseline),
    LaBraM-style patch transformer, BIOT-style FFT-tokenized transformer, and
    EEGPT-style hierarchical (spatial pool + temporal transformer).
  \item \textbf{Ablation study} (\S\ref{sec:ablation}): systematic sweep over SSL objectives
    (VICReg / SIGReg), corruption augmentations, spectral regularisation, and encoder scaling.
  \item \textbf{VICReg coefficient fix}: identified that the loss was using weights $(1,1,1)$
    instead of the paper's $(25,25,1)$; invariance was underweighted by $25\times$
    (\S\ref{sec:bugs}).
  \item \textbf{EEGPT collapse diagnosis}: channel-pool architectural bottleneck combined with
    the underweighted invariance caused EEGPT to train \emph{below} random-init performance
    (\S\ref{sec:bugs}).
  \item \textbf{Evaluation protocol audit}: caught and corrected a per-window vs per-recording
    mismatch that had inflated apparent results by $\approx$5\,pp (\S\ref{sec:eval_trap}).
  \item \textbf{Verified SOTA table} (\texttt{SOTA\_TABLE.md}): every published cell checked
    against the source PDF; two prior hallucinated values retracted.
  \item \textbf{EEGNet and ShallowConvNet reimplementations} as supervised upper bounds.
\end{enumerate}

\section{Method}

\subsection{Encoders}\label{sec:encoders}
All encoders map $x \in \mathbb{R}^{B \times 19 \times 2000}$ to
$z \in \mathbb{R}^{B \times 256}$ via \texttt{represent(x)}.

\begin{description}
  \item[\textbf{Conv}] Four strided Conv1d blocks ($k=7$, stride 2, BN + GELU),
    halving time at each step: $T=2000 \to 125$. Global average pool. 0.4\,M params.
  \item[\textbf{LaBraM}] Patchify each channel into 10 time-patches of 200 samples;
    embed with learnable channel + position embeddings; standard ViT transformer.
    $19 \times 10 = 190$ tokens. 1.2\,M params.
  \item[\textbf{BIOT}] Same tokenisation as LaBraM but embeds the FFT magnitude
    spectrum of each patch instead of raw samples. Discards phase / temporal dynamics.
  \item[\textbf{EEGPT}] Hierarchical: cross-channel attention-pool collapses
    19 channel tokens per time-patch into one summary token, then a temporal
    transformer runs over 10 summary tokens. Designed for channel-invariant features.
\end{description}

\subsection{SSL Objectives}
\textbf{VICReg} \cite{bardes2022vicreg}: two views pass through encoder + projector MLP;
loss $= \lambda\|z_1-z_2\|^2 + \mu\,V(z_1,z_2) + \nu\,C(z_1,z_2)$
with $(\lambda,\mu,\nu)=(25,25,1)$ (paper defaults).

\textbf{SIGReg} (BCS): replaces variance–covariance with an Epps–Pulley Gaussianity
test on random projections, following the LeJEPA framework.

\textbf{Corruption augmentation} (our design): on top of the standard noise/jitter/channel-drop,
we add 20\% time masking and 20\% $\pm6\sigma$ outlier injection, forcing the encoder to recover
clean structure from severely corrupted views.

\section{Evaluation Protocol}\label{sec:eval_trap}

\paragraph{The trap.}
TUAB papers (BIOT, LaBraM, FEMBA, EEGPT) report \emph{per-window (per-sample)} metrics:
every 10\,s segment is one test point with the label of its recording (409\,455 windows from
2\,339 recordings in their eval split). Our evaluation produces \emph{per-recording} scores
by mean-pooling 16 windows per patient into one embedding.

Per-recording aggregation is clinically more meaningful (you classify the \emph{patient},
not the window) but yields scores $\approx5$\,pp higher than per-window for the same encoder,
due to noise averaging. \textcolor{warn}{\textbf{Comparing per-recording to per-window numbers
overstates our results by 5\,pp.}} All comparisons below use the correct per-window metric.

\paragraph{What this means.}
Our initial claim of ``matching LaBraM'' used per-recording BACC (0.836) against
LaBraM-Base's per-window BACC (0.814). The fair per-window number for our best encoder
is \textbf{0.775}, which is below both BIOT (0.802) and LaBraM-Base (0.814).

\section{Results}\label{sec:results}

\begin{figure}[h]
  \centering
  \includegraphics[width=\linewidth]{figures/comparison_dual.png}
  \caption{Left: per-window results: the only fair comparison to the literature.
           Right: per-recording: our-side only, not comparable to published numbers.
           \textbf{Our best per-window: BACC 0.775 (SIGReg+corruption).}}
  \label{fig:comparison}
\end{figure}

\begin{table}[h]
\centering
\caption{Main results: per-window BACC / AUROC (fair comparison to literature)}
\small
\begin{tabular}{lccl}
\toprule
Method & BACC & AUROC & Note \\
\midrule
ContraWR & 0.775 & 0.846 & Literature (LaBraM Table 1) \\
BIOT (pretrained) & 0.802 & 0.874 & Literature (BIOT Table 4) \\
LaBraM-Base & 0.814 & 0.902 & Literature (LaBraM Table 2) \\
LaBraM-Huge & 0.826 & 0.916 & Literature \\
\midrule
\textbf{Ours: EB-JEPA base (frozen)} & \textbf{0.756} & 0.845 & Our run \\
\textbf{Ours: +corruption (frozen)} & \textbf{0.770} & 0.848 & Our run \\
\textbf{Ours: +spectral (frozen)} & 0.765 & 0.849 & Our run \\
\textbf{Ours: +corruption SIGReg} & \textbf{0.775} & \textbf{0.856} & Our best \\
EEGNet (supervised, ours) & 0.796 & 0.882 & Our reimplementation \\
\bottomrule
\end{tabular}
\label{tab:main}
\end{table}

\paragraph{Transfer to TUEV (6-class events).}
Frozen TUAB-pretrained encoder probed on TUEV event classification: BACC 0.364
vs random floor 0.337. The SSL representation transfers across tasks ($+$0.03 BACC,
$+$0.03 Cohen-$\kappa$), confirming it captures structure beyond recording-level label.

\section{Ablation Study}\label{sec:ablation}

\begin{figure}[h]
  \centering
  \includegraphics[width=0.85\linewidth]{figures/ablation_what_moved.png}
  \caption{Ablation over SSL objectives and encoder variants.
    Corruption augmentation is the dominant gain; scaling and masked JEPA did not help.}
  \label{fig:ablation}
\end{figure}

Key findings:
\begin{itemize}
  \item \textbf{Corruption augmentation} $+$0.014 BACC over base (per-window).
  \item \textbf{SIGReg $>$ VICReg}: $+$0.005 per-window, $+$0.009 AUROC, consistent with
    EB-JEPA paper's ``SIGReg $\geq$ VICReg'' finding.
  \item \textbf{Scaling did not help} (150 ep, deeper encoder): per-window flat at 0.768;
    the bottleneck is encoder capacity, not training time.
  \item \textbf{Masked-prediction JEPA} (EMA target + GRU predictor): BACC 0.682,
    significantly worse. A frozen global-pool readout is poorly matched to a
    frame-prediction objective; needs per-frame probing or more tuning.
  \item \textbf{Multi-corpus pretraining} (4$\times$ data): frozen 0.812 per-recording,
    fine-tuned 0.812, identical to TUAB-only fine-tuned (0.812).
    Binding constraint is encoder capacity, not data quantity.
\end{itemize}

\section{Failure Analysis}\label{sec:bugs}

\subsection{VICReg Coefficient Bug}
The implementation used $(\lambda,\mu,\nu)=(1,1,1)$ instead of $(25,25,1)$ from the paper.
The invariance term was underweighted by $25\times$ relative to covariance, so the optimizer
had little incentive to match the two views. Fixed by adding \texttt{inv\_coeff} to
\texttt{VICRegLoss} and setting $\lambda=25$ in the EEG config.

\subsection{EEGPT Encoder Collapse}
With old coefficients, EEGPT's invariance loss was stuck at $0.103$ at epoch 19
(vs Conv's $0.038$), and the trained encoder performed \emph{below} random initialisation
(accuracy 0.641 vs random floor 0.674). Two compounding causes:
\begin{enumerate}
  \item \textbf{Architectural bottleneck}: cross-channel attention-pool compresses
    19 channel tokens $\to$ 1 summary token per time-patch, giving only 10 total tokens.
    Random channel-drop augmentation applied to different channels per view
    produces very different summary tokens, making the invariance term unlearnable.
  \item \textbf{Underweighted invariance} (see above): the optimizer was not penalised
    enough to close the view gap despite the bottleneck.
\end{enumerate}
With the corrected coefficients $(25,25,1)$, EEGPT's invariance loss reaches $0.018$
by epoch 11 (vs old epoch-19 value of $0.103$); see Figure~\ref{fig:fix}.

\begin{figure}[h]
  \centering
  \includegraphics[width=0.7\linewidth]{figures/vicreg_fix.png}
  \caption{EEGPT invariance loss: old coefficients $(1,1,1)$ stuck at 0.103;
    fixed coefficients $(25,25,1)$ reach 0.018 by epoch 11.}
  \label{fig:fix}
\end{figure}

\section{Why the Per-Window Mistake Matters}

The per-recording vs per-window confusion is not merely a bookkeeping error: it reveals
a genuine tension in EEG benchmarking:

\begin{itemize}
  \item \textbf{Clinically}, the right unit is the \emph{patient/recording}
    (you want to classify patients, not 10-second snippets). Per-recording aggregation
    is arguably the \emph{correct} metric for clinical deployment.
  \item \textbf{Benchmarking}, all published SSL baselines (BIOT, LaBraM, FEMBA) use
    per-window to report higher $n$ and lower variance, making the benchmark
    look more powerful than the clinical task demands.
  \item \textbf{Our position}: per-window 0.775 (below BIOT/LaBraM); per-recording
    0.819 with a frozen probe and 0.812 with fine-tuning (both 3-seed, final epoch), a
    statistical tie with supervised EEGNet (per-recording 0.812). An earlier best-epoch
    reading of 0.837 was test-set peeking and is retracted.
\end{itemize}

\section{Conclusion}
EB-JEPA applied to EEG reaches per-window BACC 0.775 with corruption augmentation and SIGReg,
competitive with simple supervised baselines and established SSL methods like ContraWR,
but below BIOT and LaBraM at the fair comparison level.

\paragraph{What works.} Corruption augmentation. SIGReg. Strided Conv1D encoder (simple,
effective, beats all transformer variants with 25$\times$ fewer parameters).

\paragraph{What doesn't.} Masked-prediction JEPA (underperforms at 20 ep with global-pool probe).
Transformer encoders trained with wrong VICReg coefficients. Per-window scaling did not help.

\paragraph{What remains open.} EEGPT with fixed VICReg coefficients (running); neural tokenizer
(à la LaBraM) for higher-quality patch embeddings; longer pretraining on multi-corpus data with
a larger transformer encoder.

%\bibliographystyle{plain}
%\bibliography{refs}
% Manual references to keep this self-contained:
\begin{thebibliography}{9}
\bibitem{bardes2022vicreg} Bardes et al., \emph{VICReg}, ICLR 2022.
\bibitem{yang2023biot} Yang et al., \emph{BIOT}, NeurIPS 2023.
\bibitem{jiang2024labram} Jiang et al., \emph{LaBraM}, ICLR 2024.
\bibitem{zhang2024eegpt} Zhang et al., \emph{EEGPT}, 2024.
\bibitem{tang2025femba} Tang et al., \emph{FEMBA}, 2025.
\bibitem{schirrmeister2017} Schirrmeister et al., \emph{Deep Learning with CNNs for EEG}, 2017.
\end{thebibliography}

\end{document}
